\chapter{O ciclo de aprendizado de mundo aberto}
\label{cap:mundoaberto}

\begin{objetivos}
Ao final deste capítulo, você será capaz de:
\begin{itemize}[leftmargin=*,itemsep=0.2em]
  \item descrever as seis etapas do ciclo de aprendizado de mundo aberto e o papel de
        cada uma;
  \item descrever a construção de um espaço latente por autoencoder recorrente e
        justificar essa escolha frente a alternativas variacionais e de
        \ing{clustering} profundo;
  \item combinar classificadores binários e OCSVM em um voto ponderado por classe;
  \item explicar o procedimento publicado para estimar automaticamente o número de
        agrupamentos com uma floresta aleatória treinada sobre índices internos de
        qualidade;
  \item avaliar criticamente o desempenho do ciclo completo sobre uma anomalia
        genuinamente inédita.
\end{itemize}
\end{objetivos}

\section{A questão em aberto}

O \cref{cap:binarios} entregou um sistema capaz de dizer ``nenhuma das classes que
conheço''. É um avanço real e, ao mesmo tempo, uma resposta insatisfatória. Se todo
evento inédito recebe o mesmo rótulo --- \texttt{desconhecido} --- então, após alguns
meses de operação, o sistema terá acumulado milhares de janelas rejeitadas em um único
balde indistinto, sem qualquer estrutura.

O problema tem nome: \emph{rejeitar não é aprender}. Este capítulo fecha o ciclo.

\begin{definicao}{Aprendizado de mundo aberto}
Processo iterativo em que um sistema (i) reconhece o que sabe, (ii) rejeita o que não
sabe, (iii) organiza o material rejeitado em categorias candidatas, (iv) as caracteriza
e valida, e (v) as incorpora ao seu repertório, retornando a (i) com um conjunto de
classes maior.
\end{definicao}

A diferença em relação ao aprendizado incremental clássico é que aqui \emph{as classes
novas não são fornecidas por um professor}: elas emergem dos dados.

\section{As seis etapas}

\begin{enumerate}[leftmargin=*,itemsep=0.4em]
  \item \textbf{Representação.} Um autoencoder recorrente comprime a janela bruta em um
        espaço latente de baixa dimensão.
  \item \textbf{Reconhecimento.} Classificadores binários operando \emph{sobre o espaço
        latente} identificam as classes conhecidas.
  \item \textbf{Rejeição.} Um voto ponderado entre o escore binário e uma OCSVM decide
        o que é desconhecido.
  \item \textbf{Agrupamento.} O material rejeitado é particionado em grupos candidatos.
  \item \textbf{Refinamento e validação.} Os grupos são depurados por um limiar
        adaptativo e avaliados por índices internos.
  \item \textbf{Incorporação.} Grupos validados tornam-se \emph{classes virtuais}, sobre
        as quais novos classificadores binários são treinados.
\end{enumerate}

\begin{destaque}[Onde está o humano]
Em nenhuma das seis etapas o ciclo exige rotulação humana para operar. Ele exige um
humano para \emph{nomear} as classes virtuais --- e é precisamente essa lacuna que a
\cref{parte:explicabilidade} atacará, propondo que um modelo de linguagem produza a
descrição inicial que o especialista então valida ou corrige.
\end{destaque}

\section{Etapa 1: a representação latente}

\subsection{Arquitetura}

O autoencoder recebe quatro sensores amostrados em duzentos passos de tempo:

\begin{center}
\small
\begin{tabular}{lll}
\toprule
\textbf{Camada} & \textbf{Saída} & \textbf{Papel} \\
\midrule
Entrada                        & $(200, 4)$ & janela bruta \\
LSTM(32)                       & $(200, 32)$ & codificação \\
LSTM(8)                        & $(8)$       & compressão temporal \\
Densa(16)                      & $(16)$      & \textbf{espaço latente} \\
LSTM(8) $\to$ LSTM(32)         & $(200, 32)$ & decodificação \\
Densa$(4 \times 200)$          & $(200, 4)$  & reconstrução \\
\bottomrule
\end{tabular}
\end{center}

A compressão é de $800$ para $16$ valores, um fator de $50$ --- bem menos agressiva que
a do \cref{cap:dual}, porque aqui o latente não serve apenas para detectar desvio: ele
precisa \emph{discriminar} entre classes.

\subsection{Por que não um VAE ou um DCN}

A pergunta é legítima. Um autoencoder variacional impõe estrutura probabilística ao
latente; um \ing{deep clustering network} otimiza simultaneamente reconstrução e
separabilidade de agrupamentos. Ambos parecem, no papel, mais adequados à tarefa.

\begin{table}[htbp]
\centering
\caption{Acurácia do ciclo completo em função da arquitetura de representação.}
\label{tab:aberto-arquiteturas}
\begin{tabular}{lc}
\toprule
\textbf{Arquitetura da representação} & \textbf{Acurácia} \\
\midrule
Autoencoder LSTM              & $\bm{0{,}730}$ \\
Autoencoder variacional (VAE) & $0{,}727$ \\
\ing{Deep clustering network} (DCN) & $0{,}725$ \\
\bottomrule
\end{tabular}
\end{table}

\begin{resultadochave}
A diferença entre as três arquiteturas é de meio ponto percentual. O autoencoder mais
simples é o melhor dos três, e a sofisticação adicional do VAE e do DCN não se converteu
em desempenho.
\end{resultadochave}

Esse é um resultado negativo, e ele foi mantido no livro deliberadamente. Duas
interpretações são possíveis. A primeira: o gargalo do sistema não está na
representação, mas em outro ponto da cadeia --- e as etapas seguintes mostrarão que ele
está no agrupamento. A segunda: métodos que impõem estrutura ao latente (gaussianidade,
no VAE; agrupamento, no DCN) só ajudam quando essa estrutura corresponde à realidade dos
dados, e séries multivariadas de poços não são obviamente gaussianas nem obviamente
agrupáveis em esferas.

\begin{figure}[htbp]
\centering
\includegraphics[width=\textwidth]{autoencoder_clustering.png}
\caption{Projeção t-SNE do espaço latente do codificador. À esquerda, os pontos coloridos
pelos oito rótulos originais do 3W; à direita, pelos agrupamentos obtidos com $K$-médias
sobre a mesma representação. A correspondência entre as duas partições é parcial: alguns
rótulos ocupam regiões bem definidas, outros se dispersam por vários agrupamentos.
Adaptado de \citet{Lopes2025EAAI}.}
\label{fig:aberto-agrupamento-latente}
\end{figure}

A \cref{fig:aberto-agrupamento-latente} antecipa esse diagnóstico. Se a representação
fosse o gargalo, o painel da direita seria uma versão embaralhada do da esquerda. Não é:
a estrutura existe, e o agrupamento a recupera em boa parte. O que falha é a
\emph{fronteira} entre classes vizinhas --- e nenhuma variante de autoencoder resolve
isso, porque todas otimizam reconstrução, e reconstrução não distingue vizinhos.

\begin{destaque}[A releitura pelo \cref{cap:mahalanobis}]
Há uma terceira interpretação, que só ficará disponível quando a evolução narrada na
\cref{parte:evolucao} for apresentada: talvez o problema não seja \emph{qual}
autoencoder, mas o fato de usar um autoencoder. O \cref{cap:mahalanobis} mostrará que o
espaço latente extraído da penúltima camada de um classificador multiclasse é
substancialmente melhor --- silhueta de $0{,}2299$ contra $-0{,}0149$ do autoencoder.
A razão é intuitiva em retrospecto: um autoencoder é treinado para \emph{reconstruir},
não para \emph{separar}.
\end{destaque}

\section{Etapas 2 e 3: reconhecer e rejeitar}

\subsection{Classificadores binários no espaço latente}

Sobre o vetor latente, um classificador binário por classe:

\begin{center}
\small
\begin{tabular}{ll}
\toprule
\textbf{Camada} & \textbf{Configuração} \\
\midrule
Entrada        & vetor latente \\
Densa(64)      & ReLU \\
\ing{Dropout}  & $0{,}3$ \\
Densa(32)      & ReLU \\
\ing{Dropout}  & $0{,}2$ \\
Densa(1)       & sigmoide \\
\bottomrule
\end{tabular}
\end{center}

Note a diferença em relação ao \cref{cap:binarios}: lá os classificadores recebiam a
série bruta e precisavam de camadas recorrentes; aqui recebem um vetor de dezesseis
dimensões e bastam camadas densas. O trabalho de extrair estrutura temporal foi
transferido para o autoencoder --- o que é, ao mesmo tempo, uma economia e um risco: o
que o autoencoder descartou está perdido para sempre.

\subsection{O voto ponderado}

Cada classe dispõe de dois escores: o do classificador binário, $s_i^{(b)}$, e o da
OCSVM ajustada aos dados daquela classe no latente, $s_i^{(o)}$. A decisão combina os
dois:
\begin{equation}
  S_i = w_i^{(b)} \, s_i^{(b)} + w_i^{(o)} \, s_i^{(o)},
  \qquad w_i^{(b)} + w_i^{(o)} = 1,
  \label{eq:aberto-voto}
\end{equation}
com rejeição quando $\max_i S_i < \tau$, $\tau = 0{,}65$.

Os pesos foram ajustados por classe, no conjunto de validação
(\cref{tab:aberto-pesos}).

\begin{table}[htbp]
\centering
\caption{Pesos do voto ponderado por classe. $w^{(b)}$ é o peso do classificador
binário e $w^{(o)}$ o da OCSVM.}
\label{tab:aberto-pesos}
\small
\begin{tabular}{lcc}
\toprule
\textbf{Classe} & $\bm{w^{(b)}}$ & $\bm{w^{(o)}}$ \\
\midrule
1 --- Aumento de BSW          & $0{,}65$ & $0{,}35$ \\
2 --- Fechamento de DHSV      & $0{,}54$ & $0{,}46$ \\
3 --- Golfadas severas        & $0{,}75$ & $0{,}25$ \\
4 --- Instabilidade de fluxo  & $0{,}52$ & $0{,}48$ \\
5 --- Perda de produtividade  & $0{,}58$ & $0{,}42$ \\
6 --- Restrição rápida no PCK & $0{,}77$ & $0{,}23$ \\
7 --- Incrustação no PCK      & $0{,}62$ & $0{,}38$ \\
8 --- Hidrato na produção     & $0{,}53$ & $0{,}47$ \\
9 --- Hidrato no serviço      & $0{,}54$ & $0{,}46$ \\
\midrule
Nenhuma (rejeição)            & $0{,}26$ & $\bm{0{,}74}$ \\
\bottomrule
\end{tabular}
\end{table}

\begin{destaque}[A linha mais informativa da tabela]
É a última. Para decidir que uma janela \emph{não pertence a nenhuma classe conhecida}, o
peso ótimo da OCSVM é $0{,}74$ --- quase três vezes o do classificador binário. Para
todas as outras decisões, a proporção se inverte.

A interpretação é direta e generaliza para além deste sistema: \textbf{classificadores
discriminativos são bons para dizer o que algo é; modelos de suporte, treinados apenas
sobre dados de uma classe, são bons para dizer o que algo não é}. Cada família de método
tem o seu papel, e tentar que uma delas cumpra os dois é o erro que o
\cref{cap:introducao} descreveu.
\end{destaque}

\subsection{Desempenho dos classificadores binários latentes}

A \cref{tab:aberto-binarios} reporta o desempenho por classe em dois níveis de
granularidade discutidos no \cref{cap:aprendizado}.

\begin{table}[htbp]
\centering
\caption{Desempenho dos classificadores binários no espaço latente. As duas primeiras
colunas referem-se ao experimento consolidado; as duas últimas, à média e ao desvio
padrão entre execuções, por amostra e por instância.}
\label{tab:aberto-binarios}
\small
\begin{tabular}{lcccc}
\toprule
\textbf{Classe} & \textbf{ACC} & $\bm{F_1}$ & \textbf{ACC (méd.\,$\pm$\,dp)} &
\textbf{Instância (méd.\,$\pm$\,dp)} \\
\midrule
1 & $0{,}92$ & $0{,}93$ & $0{,}82 \pm 0{,}04$ & $1{,}00$ \\
2 & $0{,}90$ & $0{,}90$ & $0{,}85 \pm 0{,}03$ & $0{,}98 \pm 0{,}05$ \\
3 & $0{,}85$ & $0{,}87$ & $0{,}89 \pm 0{,}02$ & $1{,}00$ \\
4 & $0{,}69$ & $0{,}68$ & $0{,}64 \pm 0{,}03$ & $1{,}00$ \\
5 & $0{,}82$ & $0{,}83$ & $0{,}75 \pm 0{,}03$ & $0{,}76 \pm 0{,}05$ \\
6 & $0{,}94$ & $0{,}95$ & $0{,}75 \pm 0{,}13$ & $0{,}34 \pm 0{,}42$ \\
7 & $0{,}79$ & $0{,}83$ & $0{,}80 \pm 0{,}01$ & $1{,}00$ \\
8 & $0{,}90$ & $0{,}91$ & $0{,}80 \pm 0{,}06$ & $0{,}94 \pm 0{,}12$ \\
9 & $0{,}92$ & $0{,}92$ & $0{,}81 \pm 0{,}05$ & $1{,}00$ \\
\bottomrule
\end{tabular}
\end{table}

Três observações.

A classe 4 permanece a pior, com $0{,}69$ --- \emph{exatamente} o mesmo valor obtido no
\cref{cap:binarios} sobre o espaço bruto. Mudar a representação não mudou nada, o que
reforça o diagnóstico de que o problema está no rótulo.

A classe 6 exibe a patologia mais interessante da tabela: acurácia por amostra de
$0{,}75 \pm 0{,}13$, mas acurácia por instância de $0{,}34 \pm 0{,}42$. Um desvio padrão
de $0{,}42$ sobre uma média de $0{,}34$ significa que o classificador acerta a instância
inteira ou erra a instância inteira, sem meio-termo, dependendo da execução. É um
classificador instável, e a métrica por amostra o esconde.

\begin{destaque}[Por que reportar os dois níveis]
Sete das nove classes atingem acurácia por instância de $0{,}94$ ou mais, contra
acurácias por amostra entre $0{,}64$ e $0{,}89$. Reportar apenas o nível de instância
sugeriria um sistema quase perfeito; reportar apenas o nível de amostra sugeriria um
sistema medíocre. Ambos são verdadeiros e respondem a perguntas operacionais
diferentes: ``o sistema identifica o evento?'' e ``o sistema delimita o evento?''.
\end{destaque}

\section{O teste real: a anomalia de ICV}

Toda a discussão anterior seria acadêmica sem uma anomalia genuinamente inédita. Ela
existe: a falha em válvula de controle de intervalo (ICV), descrita no
\cref{cap:dominio}, que \emph{não pertence a nenhuma das nove classes do 3W} e cujos
dados são de poços reais.

\subsection{Como o ensemble reage}

\begin{figure}[htbp]
\centering
\includegraphics[width=\textwidth]{classifier_with_ICV_improved.png}
\caption{Distribuição dos escores dos nove classificadores binários quando apresentados
a janelas da anomalia de ICV --- uma classe que nenhum deles conhece. Seis dos nove
(1, 3, 4, 6, 8 e 9) produzem escore identicamente nulo. O classificador da classe 2
concentra-se em torno do limiar $0{,}5$, sem decidir. Os classificadores das classes 5 e
7 distribuem massa acima do limiar. A linha tracejada indica o limiar de decisão.
Adaptado de \citet{Lopes2025EAAI}.}
\label{fig:aberto-icv-classificadores}
\end{figure}

A \cref{fig:aberto-icv-classificadores} é uma das figuras mais informativas deste livro,
porque mostra simultaneamente o sucesso e o limite da rejeição.

O sucesso: seis dos nove classificadores respondem \emph{zero absoluto}. Não é
incerteza --- é recusa. O comportamento é exatamente o que a formulação um-contra-todos
do \cref{cap:binarios} prometia e que uma \ing{softmax} jamais poderia produzir, já que
os escores somariam um por construção.

O limite: os classificadores das classes 5 e 7 --- perda de produtividade e incrustação
--- atribuem probabilidade alta a parte das janelas. E a explicação, novamente, é
física, não algorítmica: uma ICV com falha \emph{de fato} restringe fluxo e \emph{de
fato} reduz produtividade. A assinatura sintomática é compartilhada. O classificador não
está errado no que observa; o rótulo é que não distingue causa de efeito.

\begin{resultadochave}
Do total de janelas de ICV, $\bm{48\,\%}$ foram atribuídas a alguma classe conhecida ---
majoritariamente 5 e 7 --- e $\bm{52\,\%}$ foram corretamente rejeitadas e encaminhadas
ao agrupamento.
\end{resultadochave}

Metade de uma anomalia genuinamente nova é reconhecida como nova. É simultaneamente um
resultado forte --- comparado a zero, que é o que um classificador de mundo fechado
entregaria --- e um resultado que deixa clara a distância até um sistema confiável.

\section{Etapas 4 e 5: agrupar, refinar, validar}

\subsection{Estimar o número de grupos}

Antes de agrupar é preciso decidir em quantos grupos. A solução adotada evita tanto o
arbítrio quanto a inspeção visual do cotovelo: uma \textbf{floresta aleatória} treinada
para prever $K$ a partir de oito características do próprio conjunto de dados.

\begin{center}
\small
\begin{tabular}{ll}
\toprule
\textbf{Característica} & \textbf{Descrição} \\
\midrule
$n$                    & número de amostras \\
Silhueta               & \cref{eq:silhueta} \\
Davies-Bouldin         & \cref{eq:db} \\
Calinski-Harabasz      & razão de dispersão entre e intra grupos \\
Inércia                & soma dos quadrados intra grupo \\
\texttt{avg\_k}        & estimativa por média de índices \\
\texttt{entropy\_k}    & estimativa por entropia da distribuição \\
Volume do envoltório   & hipervolume da caixa delimitadora dos dados \\
\bottomrule
\end{tabular}
\end{center}

A estimativa final é a média de duas delas:
\begin{equation}
  \hat{K} = \frac{\texttt{avg\_k} + \texttt{entropy\_k}}{2}.
  \label{eq:aberto-k}
\end{equation}

\begin{table}[htbp]
\centering
\caption{Indicadores do estimador do número de agrupamentos reportados por
\citet{Lopes2025EAAI}.}
\label{tab:aberto-rf}
\begin{tabular}{lc}
\toprule
\textbf{Métrica} & \textbf{Valor publicado} \\
\midrule
Acurácia                   & $98{,}3\,\%$ \\
Precisão                   & $0{,}98$ \\
Revocação                  & $0{,}86$ \\
Medida $F_1$               & $0{,}87$ \\
Erro médio absoluto em $K$ & $1{,}12$ \\
\bottomrule
\end{tabular}
\end{table}

\begin{destaque}[Resultado publicado e verificação independente]
A \cref{tab:aberto-rf} preserva os valores publicados e os atribui à fonte primária;
ela não é uma reprodução independente. Se ``acurácia'' significar acerto exato de $K$
e o MAE usar as mesmas predições, os dois indicadores não passam por uma verificação
aritmética: os $1{,}7\,\%$ de casos incorretos teriam erro absoluto médio condicionado de
$1{,}12/0{,}017\approx65{,}9$ grupos. Esse valor é incompatível com a escala descrita
para a tarefa.

As explicações possíveis incluem denominadores distintos, outra definição de acurácia,
MAE calculado para um componente intermediário ou erro de transcrição. Sem as predições
por conjunto ou material suplementar, não é possível escolher entre elas. Portanto, os
valores podem ser citados individualmente como resultados do artigo, mas não sustentam
a conclusão de que o estimador acerta exatamente $K$ em $98{,}3\,\%$ dos casos e erra,
em média, apenas um grupo quando falha.
\end{destaque}

\subsection{Interseção entre dois algoritmos}

O agrupamento não é confiado a um único método. Executam-se \textbf{K-means} --- com o
$K$ estimado --- e \textbf{MeanShift}, cuja largura de banda é determinada pelos próprios
dados, e retêm-se as amostras em que os dois concordam.

\begin{figure}[htbp]
\centering
\includegraphics[width=0.85\textwidth]{latent_space_with_ICV_corrected.png}
\caption{Projeção t-SNE do espaço latente. Cada cor corresponde a um agrupamento
identificado pelo K-means sobre as classes conhecidas; os círculos vermelhos são as
janelas de ICV rejeitadas, e a região sombreada é o envoltório convexo que elas ocupam.
Note que as janelas de ICV não formam um grupo compacto: elas se distribuem em subgrupos
separados, no limite entre agrupamentos existentes. Adaptado de \citet{Lopes2025EAAI}.}
\label{fig:aberto-latente}
\end{figure}

\begin{figure}[htbp]
\centering
\includegraphics[width=0.85\textwidth]{meanshift_with_ICV_corrected.png}
\caption{Agrupamento das janelas de ICV rejeitadas por MeanShift, projetado por t-SNE.
O algoritmo identifica dois grupos e marca como ruído as amostras que não pertencem a
nenhum deles --- comportamento desejável, já que a rejeição inevitavelmente arrasta
janelas de transição. Adaptado de \citet{Lopes2025EAAI}.}
\label{fig:aberto-meanshift}
\end{figure}

\begin{figure}[htbp]
\centering
\includegraphics[width=0.85\textwidth]{dbscan_with_ICV_corrected.png}
\caption{Agrupamento das mesmas janelas por DBSCAN. O algoritmo não exige que se declare
o número de grupos, mas transfere a dificuldade para os parâmetros $\varepsilon$ e
$\mathit{minPts}$: com a densidade heterogênea do espaço latente, ele ou funde os
subgrupos ou classifica a maior parte das amostras como ruído. Adaptado de
\citet{Lopes2025EAAI}.}
\label{fig:aberto-dbscan}
\end{figure}

A comparação entre as \cref{fig:aberto-latente,fig:aberto-meanshift} explica por que a
interseção é necessária. Os dois algoritmos partem de premissas diferentes --- o K-means
supõe grupos esféricos de tamanho comparável; o MeanShift, modos de densidade --- e a
concordância entre eles é um indício de estrutura genuína, não de artefato do método.
A \cref{fig:aberto-dbscan} mostra o terceiro candidato descartado: o DBSCAN é o método
teoricamente mais adequado a grupos de forma arbitrária, mas neste espaço latente sua
sensibilidade a $\varepsilon$ o torna instável demais para uso automático.

\begin{destaque}[Uma advertência sobre o t-SNE]
As \cref{fig:aberto-latente,fig:aberto-meanshift} são projeções t-SNE, e o
\cref{cap:novidades} advertiu: distâncias no t-SNE não são distâncias no latente, e o
algoritmo produz agrupamentos visualmente convincentes mesmo em ruído puro. Estas
figuras servem para \emph{ilustrar}, nunca para \emph{provar}. A evidência quantitativa
está nos índices internos e nas taxas de acerto.
\end{destaque}

\subsection{Refinamento por limiar adaptativo}

Nem toda amostra atribuída a um grupo pertence a ele. Aplica-se, dentro de cada grupo,
um limiar sobre a distância ao centroide:
\[
  T = \mu + k\sigma,
\]
com $k$ determinado pela razão $R = \mu / \sigma$ da distribuição de distâncias:
\begin{equation}
  k = \begin{cases}
    2, & \text{se } R \geq 4, \\[0.2em]
    1 + \dfrac{R-2}{2}, & \text{se } 2 \leq R < 4, \\[0.5em]
    0, & \text{se } R < 2.
  \end{cases}
  \label{eq:aberto-kadaptativo}
\end{equation}

A lógica é a seguinte. $R$ alto indica grupo compacto e homogêneo, e nesse caso pode-se
ser tolerante ($k = 2$). $R$ baixo indica grupo disperso, provavelmente contaminado, e
exige rigor ($k = 0$, ou seja, retém-se apenas o que está abaixo da distância média).
Entre os dois extremos, interpola-se linearmente.

Uma salvaguarda final: se o critério eliminar mais de dois terços do grupo, ele é
descartado e mantêm-se os dois terços mais próximos do centroide. Sem essa proteção, um
grupo mal formado poderia ser reduzido a um punhado de amostras e ainda assim promovido
a classe.

\subsection{Validação}

Grupos sobreviventes são avaliados pelos índices internos do \cref{cap:novidades}:
silhueta, Davies-Bouldin e Calinski-Harabasz. Um grupo só é promovido a classe virtual
se apresentar coesão e separação adequadas segundo os três.

\subsection{O que o agrupamento acerta, quando acerta}

A acurácia agregada esconde uma variabilidade grande entre combinações de classes
desconhecidas. Vale examinar um caso favorável para saber o que o método produz quando
as condições lhe são proprícias.

\begin{figure}[htbp]
\centering
\begin{minipage}{0.49\textwidth}
  \centering
  \includegraphics[width=\textwidth]{best_original.png}
\end{minipage}\hfill
\begin{minipage}{0.49\textwidth}
  \centering
  \includegraphics[width=\textwidth]{best_clustering.png}
\end{minipage}
\caption{Ocultando as classes 3, 6 e 8 e tratando-as como desconhecidas. À esquerda, a
distribuição original dos dados, colorida pelos rótulos verdadeiros; à direita, os três
agrupamentos recuperados por $K$-médias sem acesso a esses rótulos. Adaptado de
\citet{Lopes2025EAAI}.}
\label{fig:aberto-melhor-caso}
\end{figure}

A \cref{fig:aberto-melhor-caso} mostra esse caso: escondem-se três das nove classes
conhecidas, treina-se o sistema com as seis restantes e verifica-se o que o agrupamento
faz com o material rejeitado. A correspondência entre os dois painéis é boa --- medida
$F_1$ de $0{,}71$, $0{,}97$ e $0{,}76$ para as classes 3, 6 e 8 --- e o número de grupos
foi estimado corretamente.

\begin{figure}[htbp]
\centering
\includegraphics[width=0.78\textwidth]{confusion_matrix_kmeans.png}
\caption{Matriz de confusão normalizada do agrupamento por $K$-médias, agregada sobre as
$\binom{9}{3} = 168$ combinações possíveis de três classes ocultas. A diagonal domina,
mas os erros não estão distribuídos de modo uniforme. Adaptado de
\citet{Lopes2025EAAI}.}
\label{fig:aberto-matriz-kmeans}
\end{figure}

Um caso favorável, contudo, não é evidência. Por isso o experimento foi repetido para
\emph{todas} as $168$ combinações de três classes escolhidas entre nove, e a
\cref{fig:aberto-matriz-kmeans} agrega o resultado. A leitura importante não é a
diagonal, e sim o padrão fora dela: os erros concentram-se nos mesmos pares que já
haviam aparecido na \cref{fig:aberto-latente} e na rejeição da anomalia de ICV. Não é
ruído do algoritmo --- é sobreposição real de assinatura entre classes vizinhas, e nenhum
método de agrupamento a desfaz sozinho.

\section{Etapa 6: classes virtuais}

\begin{definicao}{Classe virtual}
Agrupamento validado que passa a ser tratado como classe pelo sistema, recebendo seu
próprio classificador binário, mas que ainda não foi nomeado nem confirmado por um
especialista.
\end{definicao}

A promoção a classe virtual é a etapa que fecha o ciclo, e ela é barata precisamente por
causa da propriedade de extensibilidade incremental estabelecida no \cref{cap:binarios}:
treina-se um novo classificador binário, e nenhum dos existentes é tocado.

\begin{figure}[htbp]
\centering
\includegraphics[width=\textwidth]{classifier_with_ICV_with_10_classes.png}
\caption{O mesmo cenário da \cref{fig:aberto-icv-classificadores}, agora com dez
classificadores. O décimo foi treinado exclusivamente sobre as janelas antes rejeitadas
e reivindica a maior parte delas; as reivindicações espúrias das classes 2 e 5
retrocedem. Adaptado de \citet{Lopes2025EAAI}.}
\label{fig:aberto-dez-classes}
\end{figure}

A \cref{fig:aberto-dez-classes} é o fecho experimental do ciclo, e deve ser lida em par
com a \cref{fig:aberto-icv-classificadores}. São os mesmos dados de ICV apresentados ao
mesmo conjunto de classificadores, com uma única diferença: um décimo classificador foi
acrescentado, treinado apenas com o material que o próprio sistema havia separado como
desconhecido. O que era detecção dispersa entre duas classes erradas torna-se uma
reivindicação concentrada e correta --- a validação de que $0{,}98$ de acurácia na
décima classe não custou desempenho nas nove anteriores.

\begin{destaque}[O que esta figura não prova]
O décimo classificador aprendeu a reconhecer a falha de ICV, mas ninguém lhe disse que
aquilo \emph{era} uma falha de ICV. Para o sistema, a classe continua anônima: um
agrupamento coerente, reprodutível e agora reconhecível, sem nome nem causa.

Nomear é outro problema, e é exatamente onde começa a \cref{parte:explicabilidade}.
\end{destaque}

\begin{resultadochave}
Sobre o conjunto de janelas rejeitadas, o ciclo de agrupamento e refinamento atinge
acurácia de $\bm{80{,}8\,\%}$ no nível de amostra na atribuição a grupos coerentes.
\end{resultadochave}

\section{Balanço}

O ciclo funciona, e funciona parcialmente. Vale enunciar as duas metades com igual
clareza.

\textbf{O que funciona.} A rejeição não é uma promessa teórica: $52\,\%$ de uma anomalia
genuinamente inédita foram identificados como tal, com seis de nove classificadores
respondendo zero absoluto. O agrupamento subsequente organiza esse material com
$80{,}8\,\%$ de acurácia. A estimativa automática de $K$ dispensa arbítrio. E o
mecanismo de classes virtuais permite crescimento do repertório sem retreinamento
global.

\textbf{O que não funciona.} Quase metade das janelas de ICV foi absorvida pelas classes
sintomáticas 5 e 7. A escolha da representação --- autoencoder, VAE ou DCN --- é
irrelevante, o que sugere que o gargalo não foi identificado. As classes 4 e 5 falham
sistematicamente. E, sobretudo, \textbf{as classes virtuais não têm nome}: o sistema
produz ``grupo 11'', não ``falha de ICV''.

\begin{destaque}[As três portas que este capítulo abre]
Cada uma dessas limitações define uma das partes seguintes do livro.

A representação inadequada leva ao \cref{cap:mahalanobis}: o latente do classificador
multiclasse supera o do autoencoder por uma margem que torna o resultado da
\cref{tab:aberto-arquiteturas} compreensível em retrospecto.

A detecção grosseira, no nível de janela inteira, leva ao \cref{cap:segmentacao}: uma
U-Net que segmenta ponto a ponto, elevando a acurácia de $0{,}500$ para $0{,}863$.

E a ausência de nomes leva à \cref{parte:explicabilidade}, onde um modelo de linguagem
recebe as estatísticas do segmento anômalo e produz uma descrição em linguagem natural
--- que, para a maioria das classes, converge para um nome consistente.
\end{destaque}

\resumocapitulo{%
O ciclo de aprendizado de mundo aberto tem seis etapas: representação, reconhecimento,
rejeição, agrupamento, refinamento e incorporação. A representação é um autoencoder LSTM
que comprime $200 \times 4$ valores em dezesseis dimensões; alternativas variacionais e
de \ing{clustering} profundo produzem desempenho estatisticamente indistinguível
($0{,}730$, $0{,}727$ e $0{,}725$), o que indica que o gargalo está em outro lugar. A
rejeição combina classificadores binários e OCSVM em um voto ponderado por classe, e o
peso ótimo da OCSVM na decisão de rejeitar é $0{,}74$ contra $0{,}26$ do classificador
--- evidência de que reconhecer e rejeitar são tarefas para famílias distintas de
método. Sobre uma anomalia genuinamente inédita, a falha de ICV, seis dos nove
classificadores produzem escore nulo e $52\,\%$ das janelas são corretamente rejeitadas;
os $48\,\%$ restantes são absorvidos pelas classes sintomáticas 5 e 7. O agrupamento usa
a interseção entre K-means e MeanShift, com $K$ estimado por uma floresta aleatória
sobre oito índices. O artigo reporta $98{,}3\,\%$ de acurácia e MAE de $1{,}12$, valores
preservados aqui com a ressalva de que suas definições conjuntas aguardam reconciliação
a partir das predições por conjunto. O refinamento usa um limiar adaptativo
$\mu + k\sigma$ com $k$ dependente da razão
$R = \mu/\sigma$. A acurácia do ciclo de agrupamento é de $80{,}8\,\%$ no nível de
amostra. Os grupos validados tornam-se classes virtuais --- que ainda não têm nome.}

\begin{exercicios}
\item O peso ótimo da OCSVM na decisão de rejeição é $0{,}74$, contra valores entre
      $0{,}23$ e $0{,}48$ nas decisões de reconhecimento. Explique essa assimetria em
      termos da natureza discriminativa dos classificadores binários e da natureza
      generativa dos modelos de suporte.

\item Implemente o limiar adaptativo da \cref{eq:aberto-kadaptativo}. Gere conjuntos
      sintéticos com diferentes graus de dispersão e verifique se o comportamento de $k$
      corresponde à intenção declarada. Proponha e avalie uma alternativa contínua e
      diferenciável.

\item Mostre que, se a acurácia de acerto exato de $K$ é $a$ e o erro médio absoluto é
      $m$, então o erro absoluto médio condicionado aos casos incorretos é
      $m/(1-a)$. Use essa identidade para propor verificações de consistência para um
      relatório que apresente as duas métricas.

\item \textbf{Projeto de extensão.} Com acesso aos dados e a um protocolo explicitamente
      versionado, substitua o autoencoder por um \ing{transformer} de séries temporais e
      compare-o com a \cref{tab:aberto-arquiteturas}. As diferenças entre representações
      permanecem pequenas? Documente por que seus valores podem divergir dos publicados e
      o que o resultado implica sobre onde concentrar esforço de engenharia.

\item A interseção entre K-means e MeanShift descarta amostras em que os dois algoritmos
      discordam. Quantifique, sobre o 3W, a fração descartada. Discuta o compromisso
      entre pureza dos grupos e cobertura, e proponha um critério para decidir quando a
      interseção é preferível à união.

\item Argumenta-se que $48\,\%$ das janelas de ICV foram absorvidas pelas classes 5 e 7
      por compartilharem assinatura sintomática. Projete um experimento que teste essa
      hipótese contra a alternativa de que se trata de simples limitação do
      classificador.

\item Uma classe virtual é promovida sem nome nem confirmação humana. Discuta os riscos
      operacionais de acumular classes virtuais indefinidamente e proponha uma política
      de governança: quando promover, quando fundir, quando descartar, e com que
      frequência solicitar revisão humana.
\end{exercicios}

\begin{leituras}
\item \citet{LopesTese2025} --- a tese que originou este capítulo, com detalhamento
      completo do ciclo.
\item \citet{Lopes2025EAAI} --- o artigo em que o ciclo foi publicado, origem das
      figuras deste capítulo.
\item \citet{openworld1}, \citet{openworld2} e \citet{openworld3} --- os fundamentos do
      aprendizado de mundo aberto.
\item \citet{Pimentel2014review} --- revisão abrangente de detecção de novidade.
\item \citet{Breiman2001} --- florestas aleatórias, aqui usadas de forma incomum, como
      estimador de metaparâmetro.
\item \citet{maaten2008visualizing} --- o t-SNE, com a discussão original sobre o que
      suas projeções podem e não podem sustentar.
\end{leituras}
